\documentclass[10pt]{article}
\usepackage[margin=12mm]{geometry}
\usepackage{graphicx}
\usepackage{caption}
\pagestyle{empty}
\captionsetup{font=small,labelfont=bf}

\begin{document}
\begin{figure}[p]
\centering
\includegraphics[width=\textwidth]{figS2_ce_cell_type_cost_gain_panel.pdf}

\caption{\textbf{Fate-specific distance changes along the global terminal-cell
Pareto front.}
%
Per-cell distance changes at the travel-minimizing assignment (teal),
maximum-retention compromise (blue diamond), and cell-state-minimizing
assignment (orange) on the global Pareto front, all relative to the natural lineage and
reported in embryo-wide first-cousin-null standard-deviation units per cell.
Negative values are distance reductions and positive values are increases.
Objective-colored shading marks reductions and grey shading marks increases;
grey lines join the two single-objective endpoints, and marker area increases
with cell count. Both panels are oriented from the travel optimum at left to
the cell-state optimum at right, so the reduction and increase sides reverse
between the two objectives. Travel and cell-state changes are shown on
separate axes and interpreted independently. At the cell-state optimum,
programmed-death cells show the largest per-cell cell-state reduction and
epithelial cells the largest per-cell travel increase; at the travel optimum,
muscle cells show the largest per-cell cell-state increase. Neurons are not
the most extreme per cell but contribute strongly embryo-wide through their
abundance.}
\label{fig:ce_cell_type_cost_gain}
\end{figure}
\end{document}
